% P. Svidler - M. Carlsen, Norway Chess 2014: your comments. Format: games/README.md.
@ title | P. Svidler - M. Carlsen, Norway Chess 2014

@ 3...f5
\mythoughts{Black is playing the Grand Prix Attack.\index{Openings!Grand Prix Attack} The idea is f5 and \B b4xc3.
He castles kingside, brings the queen to h5 via e8, and the knight to g4.
That yields three attackers plus a pawn against the king. An f4 break could allow
a rook sacrifice on f3 and mate on h2.

Of course this is only a sketch; the same idea can crop up in other openings.
White could adopt a Rossolimo against the Sicilian (1.e4 c5 2.\N f3 \N c6 3.\B b5)
and mirror the whole plan with reversed colors.}

@ 7.bxc3 | board
\mythoughts{Black has already traded the knight on c3 for the bishop, as planned, and still intends to
castle kingside.

Where should the light-squared bishop go? A fianchetto is wrong: b7 is loose, and on a6
the bishop has too little to do.

I do not want to let White play \B a3 and spoil my castling, so I castle now.}

@ 7...d6
\myreflection{7...O-O allows 8.c5, leaving a weakness on d6 for Black.
My thoughts above are not entirely wrong, but as Lasker put it, ``if you find
a good move, look for a better one.''\index{Practical play!Look for a better move} 7...d6 is better.}

@ 8.O-O
\mythoughts{As planned, I castle on the kingside.}

@ 9.Rb1
\mythoughts{I bring the queen to e8 as planned.}

@ 9...Qe8
\myreflection{Magnus has followed the script I expected. Over the board, play is not only
move by move but plan by plan: a plan is a short sequence of moves.\index{Practical play!Plan by plan}
As long as nothing forces a change of course, following the plan saves time.}

@ 10.Qb3
\mythoughts{I dislike having my queen on the same diagonal as White's queen, so I would
play 10...\K h8.}

@ 10...b6
\myreflection{10...\K h8 is playable. 11.c5+ \B e6 White loses a pawn.
Judging how much prophylaxis you need is hard.\index{Strategy!Prophylaxis} If the king stays on g8, Black must watch
the c5 break for the rest of the game. Former world champion Petrosian would tuck the king on h8 without hesitation.

Peter Leko is often on commentary. He likes to say ``I don't want adventure. I don't want
to calculate,'' while a younger co-commentator rattles off concrete lines to show that all is well.

Am I getting old?}

@ 11.Nh4
\mythoughts{In closed systems where White has a light-squared fianchetto, Black often struggles to
develop on the queenside because of the pin.\index{Pin} With White's queen on b3, I should use the chance
to bring out the knight and rook.}

@ 12.Qa3
\mythoughts{We stick to the plan and develop the rook to a safer square.}

@ 13.Be3 | board
\mythoughts{Black is aiming for ...f4.
White has no time for 14.\Q xa5 bxa5 15.\R xb8. If the bishop were not
under attack, he would stand a fair chance of meeting Black's attack.

After 14.gxf4 exf4 15.\B xf4 \Q xe2 Black has some extra material
even though the attack has fizzled.}

@ 14.gxf4 | board
\mythoughts{When I calculated a move ago, I thought capturing on f4 was my only try. Now it feels weak:
Black no longer has a real attack.

14...\Q h5 is both forcing and in line with my plan to sharpen the game.}

\myreflection{My hesitation is obvious. One move earlier, capturing on f4 looked like my only option;
suddenly it looks poor, even though I already sensed that the attack was gone and only
concrete gains remained. In-between moves such as 14...\Q h5 are easy to
overlook when you calculate several moves deep.\index{Practical play!Intermediate move} They are also forcing and make the knight
retreat. As a practical rule, leading with such moves is useful: the opponent must answer,\index{Practical play!Forcing moves first}
and the tree stays smaller. I then see how strong 10...b6 was: White's play
stalls. Queen and rook are tied to the queenside; while they lack a breakthrough there,
Black keeps realistic chances on the kingside. (Sometimes White gets rolling too, but Black
may still win by mating first.)}

The computer prefers 14...e4. After 15.dxe4 \N xe4
White must react because a fork wins the exchange.\index{Knight!Fork} Black can then snap the c4-pawn. Compared
with my original scheme where the queen crashes through on e2, this leaves Black better placed:
in my line White gets an open e-file and ideas against the king. Side-by-side comparisons like
this are a handy way to choose at the board.\index{Practical play!Comparing positions}

@ 15.Nf3
\mythoughts{15...\B h3 trades a defender while drawing the queen closer to the enemy king.}

@ 17.Kh1 | board
\mythoughts{The original plan was 17...\N g4. After 18.\R g1 Black
cannot continue the attack: White has 19.\R xg4 with counterplay, and the g-file opens.
18...exf4 19.\B d4 and suddenly White eyes the black king, though before Black's 17th move
White had almost nothing active.

I also dislike 17...e4 because it hands the bishop the d4-square
against the g7-pawn.

I still need more weight on the kingside. The f8-rook is already part of the attack.
White cannot grab on e5 with the f-pawn because ...\R xf3 and mate on h2 hangs in the air.
The a8-rook is passive; ...\R be8 intends ...\R g6 or ...\R h6 via e6.}

\myreflection{The line of thought above is reasonable. The engine is not too harsh on 17...e4 because after
18.dxe4 \N xe4 19.\R g1 c5 Black is still alright. Part of me likes not forcing the attack
if a won ending is on offer; on the other hand I dislike giving White the g-file.}

@ 18.Qb2 | board
\mythoughts{18...\R e6 loses because of 19.\N g5.
18...\N g4 19.\B d2 e4 Black destroys White's pawn structure.
18...e4 needs to be calculated. 19.dxe4 \R xe4 White
cannot play 20.\R g1 because 20...\R xe3 and then mate with 21...\N e4.

I cannot tell which continuation is stronger.}

\myreflection{After 18...\N g4 19.\R g1 Black still keeps an edge. I would feel jarred:
the attack is over. Pushing ...e4 a tempo later was a sound idea, but hoping for White's 19.\B d2 was naive.}

\myreflection{I get lost calculating 18...e4 19.dxe4 because recapturing with the knight promises a clearly better ending:
19.dxe4 \N xe4 20.\R g1 \N xc4 21.\Q b3 \Q e6 22.\R g2 \N xe3 23.fxe3 \Q xb3 24.\R xb3. That sequence never appears on the board and Black's edge evaporates.}

\myreflection{The practical choice needs little calculation: 18...e4 is better because it forces the pace.
After 18...\N g4 I may easily overlook something and rebuild the plan from scratch.
The whole point of 18...\N g4 was to follow with 19...e4, so play the pawn break first.}

@ 19.Ng5
\mythoughts{White's answer is natural. After 19...\Q h5 20.dxe4 \N g4 21.\N f3 \R xe4
Black threatens ...\R f4-f3 and mate on h2. I go for it.}

\myreflection{We will see that 21...\R xe4 is not forced---Black has an even stronger option.}

@ 21.Nf3
On move 19 I missed 21...\N xc4.

@ 21...Nxc4
\myreflection{Here is a pattern I know too well: I find a good plan and a good move yet miss
a stronger one. Better moves are often easier to spot when they are forcing: the opponent must react.\index{Practical play!Forcing moves first}\index{Practical play!Look for a better move}
Grabbing on c4 is superior because the white queen must bail out at once. 21...\R xe4 22.\Q b5
Black probably still keeps an edge.}

@ 22...Rxe4
Capturing on e4 was part of the plan as well.

@ 23.Rg1 | board
\mythoughts{23...\R exf4 24.\B xf4 \R xf4 25.\R xg4 is too committal: I see no way to press for a win once everything trades down.

23...\K h8 threatens ...\N xe3. 24.\R xg4 \Q xg4 25.\Q c2 leaves h7 very tender.}

\myreflection{I failed to find the right move. 23...\K h8 is adequate, and I overlooked the game continuation.
23...\K h8 24.\R xg4 \Q xg4 25.\Q c2 \N xe3 26.fxe3 \R xe3 wins for Black.
23...\K h8 24.\R xg4 \Q xg4 25.\R g1 leaves Black up material with very active pieces and realistic winning chances.}

@ 24.Qb5 | board
\mythoughts{Black cannot improve further without landing a decisive blow---and the pieces call for one.

24...\R exf4 25.\B xf4 \R xf4 26.\R xg4 \R xg4 I am unsure Black can convert; he is only a pawn up.

24...\R xe3 25.\R xg4 \Q xg4 26.\R g1 \Q xf4 27.fxe3 \N xe3 Again Black is a pawn ahead, but White retains some activity on the g-file.

24...\N xf2+ 25.\B xf2 \R xe2 Black is crashing through.}

\myreflection{I believe I found something stronger than Carlsen's choice, yet my analysis is shaky.
24...\R exf4 25.\B xf4 \R xf4 26.\Q d7 wins a pawn, but with queens on the board White keeps some counterplay.

24...\R xe3 25.\R xg4 \Q xg4 26.\R g1 \Q xf4 27.\Q xd5+ \Q f7 28.\Q xf7+ \R xf7 29.fxe3 \N xe3 After the queens come off, Black is simply winning.\index{Strategy!Exchanging pieces when ahead}}

@ 27.Kf1
What follows is no longer my chosen continuation---I had already found a win earlier. Along the actual moves Black cannot convert, for example:

27...\Q xf3 28.\R g3 \Q e4 29.\Q d7 \Q xb1+ 30.\K xe2 \Q d1+ 31.\K xf2 \Q c2+ 32.\K g1 \Q d1+ 33.\K f2 \Q c2+ 34.\K g1

@ 28.Rxg7+
Capturing the rook also leads to a draw, for example
28...\K xg7 29.\Q d7+ \Q f7 30.\Q g4+ \Q g6 31.\Q d7+ \K h8 32.\Q d8+ \Q g8 33.\Q xg8+ \K xg8 34.\K xe2 \N xc3+ 35.\K d3 \N xb1

@ 37...Kf8
\myreflection{I am glad my calculations were not wholly off.
The Grand Prix Attack is worth keeping in mind: the ideas are fairly straightforward. Svidler was awful in the opening and middlegame and escaped only by a miracle.

I take away that forcing moves deserve the first look.\index{Practical play!Forcing moves first} I will miss things as a human, and mimicking an engine by trying to calculate everything is pointless; human-style chess is more practical.

When the position calls for direct play, the right move is often not hard to find. I still need study in these typical setups---chess offers plenty of them.

I have been watching live games from the Candidates. The human commentators explain plans in plain language without drowning the audience in lines; I should annotate the same way.}

